\documentclass{report}
\usepackage{amsmath,amssymb}
\setlength{\parindent}{0mm}
\setlength{\parskip}{1em}
\begin{document}
\begin{center}
\rule{6in}{1pt} \
{\large Chris Hansen \\
{\bf GEOMETRIC MULTIGRID FOR MHD SIMULATIONS WITH NEDELEC FINITE ELEMENTS ON TETRAHEDRAL GRIDS}}

Aeronautics and Astronautics \\ University of Washington \\ Box 352250 \\ AERB Room 120 \\ Seattle \\ WA 98195
\\
{\tt hansec@uw.edu}\\
George Marklin\end{center}

The Magneto-HydroDynamic (MHD) model is used extensively to simulate
macroscopic plasma dynamics in Magnetic Confinement Fusion (MCF) devices.
In these simulations, the span of time scales from fast wave dynamics to
the desired evolution of equilibrium due to transport processes is large,
resulting in stiff linear systems for implicit time advance. Many
existing codes leverage toroidal symmetry, a common feature of many MCF
devices, to develop efficient preconditioners for the linearized system.
Recent interest in 3D effects in symmetric MCF devices and experiments
with non-symmetric boundaries has renewed interest in improved geometric
flexibility within numerical tools.

The PSI-TET code is a 3D non-linear extended MHD code based on a high
order finite element method using unstructured tetrahedral grids,
allowing arbitrary geometry to be captured easily. A mixed FE basis is
used with a Nedelec vector basis set used to represent magnetic flux and
scalar Lagrange elements used for each component of the plasma velocity.
This talk discusses the application of geometric multigrid to
preconditioning both equilibrium and time dependent MHD solves in
PSI-TET. For equilibrium solves, the relevant linear systems are
symmetric so a multigrid-preconditioned CG with Jacobi smoothers is used.
The linearized systems produced for time dependent solves are not
symmetric so a multigrid-preconditioner GMRES method is employed, with
block-preconditioned smoothers. We discuss the implementation and
performance of these methods in comparison to single level
preconditioners for benchmark and experimental validation simulations. We
also present implementation considerations for constructing multigrid
levels from both h and p refinement of a base mesh while maintaining
consistency CAD geometry. Work supported by DOE.


\end{document}
